\section{Completed residue fields of arbitrary seminorm points}
\label{sec:general-residue-fields}

Let $A$ be an arbitrary commutative ring and let $x\in\mathcal S(A)$.
Write $\mathfrak p_x=\{a:|a|_x=0\}$. This is a prime ideal. The
constructions below use this actual support, with no field structure or
trivial valuation assumed on the coefficient ring.

\begin{definition}[Residue field and induced seminorm]
\label{spectrum:residue-field}
\lean{EnvelopeProperty.Berkovich.Spectrum.ResidueField,
EnvelopeProperty.Berkovich.Spectrum.residueValuation,
EnvelopeProperty.Berkovich.Spectrum.residueValuation_algebraMap,
EnvelopeProperty.Berkovich.Spectrum.residueAbs,
EnvelopeProperty.Berkovich.Spectrum.isNonarchimedean_residueAbs}
\leanok
\uses{spectrum:support}
\source{fhx-envelope-20261028:subsec:prelim-berkovich}
Set $\kappa(x)=\operatorname{Frac}(A/\mathfrak p_x)$.
There is a multiplicative ultrametric absolute value on $\kappa(x)$
extending $|\cdot|_x$ on $A$.
\end{definition}
\begin{proof}
The seminorm descends to $A/\mathfrak p_x$ because it vanishes on
$\mathfrak p_x$. Its support there is zero. Every nonzero denominator
therefore has positive seminorm, so the valuation extends to the fraction
field by $|a/b|=|a|/|b|$. This is a valuation on a field, hence vanishes
only at zero. Regard its nonnegative real values as a real absolute value;
the ultrametric inequality implies the usual triangle inequality and also
retains the stronger ultrametric inequality. The quotient and localization
maps preserve the seminorm of every original element of $A$.
\end{proof}

\begin{definition}[Completed residue field]
\label{spectrum:completed-residue-field}
\lean{EnvelopeProperty.Berkovich.Spectrum.H,
EnvelopeProperty.Berkovich.Spectrum.residueToH,
EnvelopeProperty.Berkovich.Spectrum.norm_residueToH,
EnvelopeProperty.Berkovich.Spectrum.denseRange_residueToH,
EnvelopeProperty.Berkovich.Spectrum.toH,
EnvelopeProperty.Berkovich.Spectrum.norm_toH,
EnvelopeProperty.Berkovich.Spectrum.ker_toH}
\leanok
\uses{spectrum:residue-field}
\source{fhx-envelope-20261028:lem:finite-flat-norm}
The completed residue field $\mathcal H(x)$ is the completion of
$\kappa(x)$ in the induced absolute value. It is a complete nonarchimedean
normed field. The canonical ring homomorphism $\iota_x:A\to\mathcal H(x)$
satisfies $\|\iota_x(a)\|=|a|_x$ and has kernel $\mathfrak p_x$.
The canonical embedding $\kappa(x)\to\mathcal H(x)$ preserves absolute
values and has dense image.
\end{definition}
\begin{proof}
Realize $\kappa(x)$ as a normed field using its absolute value and take
its uniform completion. The completion of a normed field is a complete
normed field. To retain the ultrametric inequality, observe that the
condition $\|u+v\|\le\max\{\|u\|,\|v\|\}$ is closed in the pair $(u,v)$;
it holds on the dense image of $\kappa(x)^2$, so holds on the completion.
The canonical completion map preserves the norm and has dense image.
Composing it with $A\to\kappa(x)$ gives $\iota_x$ and its norm identity.
An element maps to zero exactly when its seminorm is zero, yielding the
kernel assertion.
\end{proof}

\begin{theorem}[Residue fields under arbitrary ring pullback]
\label{spectrum:residue-field-pullback}
\lean{EnvelopeProperty.Berkovich.Spectrum.ResidueField.map,
EnvelopeProperty.Berkovich.Spectrum.ResidueField.map_algebraMap,
EnvelopeProperty.Berkovich.Spectrum.residueValuation_map,
EnvelopeProperty.Berkovich.Spectrum.ResidueField.map_id,
EnvelopeProperty.Berkovich.Spectrum.ResidueField.map_comp}
\leanok
\uses{spectrum:residue-field,spectrum:pullback}
Let $\varphi:A\to B$ be any unital ring homomorphism, $y\in\mathcal S(B)$,
and $x=\varphi^*y$. There is a canonical field embedding
$\kappa(x)\to\kappa(y)$ that preserves the induced absolute values and
commutes with the maps from $A$ and $B$. These embeddings respect identity
and composition of ring homomorphisms.
\end{theorem}
\begin{proof}
The support satisfies $\mathfrak p_x=\varphi^{-1}(\mathfrak p_y)$.
Thus $\varphi$ induces an injective homomorphism of the quotient domains
and hence a homomorphism of their fraction fields. The induced valuation
on $\kappa(x)$ obtained by pulling back that of $\kappa(y)$ agrees with
the valuation of $x$ on $A$. Valuations on a fraction field are determined
by their values on the quotient domain, so the two valuations coincide.
For identity and composition, the proposed field maps agree on the images
of all elements of $A$. The universal property of quotient and fraction
fields gives equality of the entire maps.
\end{proof}

\begin{theorem}[Completed residue field functoriality]
\label{spectrum:completed-residue-field-pullback}
\lean{EnvelopeProperty.Berkovich.Spectrum.H.map,
EnvelopeProperty.Berkovich.Spectrum.H.norm_map,
EnvelopeProperty.Berkovich.Spectrum.H.isometry_map,
EnvelopeProperty.Berkovich.Spectrum.H.map_toH,
EnvelopeProperty.Berkovich.Spectrum.H.map_id,
EnvelopeProperty.Berkovich.Spectrum.H.map_comp}
\leanok
\uses{spectrum:completed-residue-field,spectrum:residue-field-pullback}
The residue-field embedding extends canonically to an isometric ring
embedding $\mathcal H(x)\to\mathcal H(y)$, commuting with $\iota_x$,
$\iota_y$, and $\varphi$. The completed embeddings respect identity and
composition.
\end{theorem}
\begin{proof}
The residue-field map preserves norms, so it is an isometry and hence
uniformly continuous. Its completion is a ring homomorphism extending
the residue-field map. Norm preservation passes from the dense residue
field to the completion because norm is continuous. Consequently the
completed map is an isometry and is injective.
The commuting square follows first for residue-field elements, and thus
in particular for the original elements of $A$.
The identity and composition laws hold on the dense residue-field image
by the preceding theorem. Continuous maps into the Hausdorff completed
field that agree on a dense set agree everywhere, proving these laws.
\end{proof}

\begin{corollary}[Compatibility with a prescribed base seminorm]
\label{spectrum:relative-completed-residue-field}
\lean{EnvelopeProperty.Berkovich.RelativeSpectrum.H,
EnvelopeProperty.Berkovich.RelativeSpectrum.norm_algebraMap_H,
EnvelopeProperty.Berkovich.RelativeSpectrum.H.map}
\leanok
\uses{spectrum:completed-residue-field-pullback,spectrum:relative}
Let $R$ be any commutative ring with a prescribed seminorm $v$, and let
$A,B$ be commutative $R$-algebras. For $x\in\mathcal S_v(A)$, the field
$\mathcal H(x)$ is an $R$-algebra with
$\|r\|_{\mathcal H(x)}=v(r)$ for every $r\in R$.
If $\varphi:A\to B$ is an $R$-algebra homomorphism and $x=\varphi^*y$,
the completed residue-field embedding is an $R$-algebra homomorphism.
\end{corollary}
\begin{proof}
Compose $R\to A$ with $\iota_x$ to obtain the $R$-algebra structure.
The norm identity for $\iota_x$ and the defining restriction of a relative
point give the displayed base norm. The commuting square of completed
residue fields, applied to the image of $r\in R$, shows that pullback
embeddings commute with this base action, since $\varphi$ is an
$R$-algebra homomorphism.
\end{proof}

\begin{remark}\label{spectrum:completed-residue-field-scope}
\uses{spectrum:completed-residue-field,spectrum:relative-completed-residue-field}
No boundedness, finite-generation, field, or completeness assumption is
made on $A$ or the prescribed coefficient ring. The base map may have
kernel $\operatorname{supp}(v)$, so it is norm preserving in the displayed
sense without being asserted injective. The field maps under ring pullback
are isometric embeddings. These constructions do not assert an extension
theorem for valuations or a description of finite analytic fibres by
Artinian tensor factors; those are separate results.
\end{remark}
